% === Communications in Mathematical Physics === 官方投稿模板（YangTex 内置）
% 官方类：svjour3.cls | 引擎：pdflatex | 分类：Springer 数学期刊
% 官方类缺失时自动降级为 article 预览排版，仍可正常编译与双向逐行同步。
\IfFileExists{svjour3.cls}{%
  \documentclass[smallextended]{svjour3}
}{%
  \documentclass[11pt,a4paper]{article}
  \def\YTfallback{1}
}
\usepackage{amsmath,amssymb,amsthm,mathtools}
\usepackage{graphicx}
\providecommand{\journal}[1]{}
\providecommand{\journalname}[1]{}
\providecommand{\institute}[1]{}
\providecommand{\headers}[2]{}
\providecommand{\shorttitle}[1]{}
\providecommand{\subjclass}[2][]{}
\providecommand{\IEEEmembership}[1]{}
\providecommand{\affiliation}[1]{}
\providecommand{\address}[1]{}
\providecommand{\email}[1]{}
\providecommand{\smartqed}{}
\providecommand{\at}{, }
\providecommand{\eps}{\varepsilon}
\newtheorem{theorem}{Theorem}[section]
\newtheorem{lemma}[theorem]{Lemma}
\newcommand{\R}{\mathbb{R}}
\newcommand{\Om}{\Omega}
\newcommand{\dd}{\,\mathrm{d}}
\journalname{Communications in Mathematical Physics}
\title{A Large Deviation Principle for Stochastic Navier-Stokes Equations}
\author{Yang Jianwei}
\date{}
\begin{document}
\maketitle
\ifdefined\YTfallback
\noindent\fbox{\parbox{0.96\linewidth}{\small\textbf{[YangTex preview mode] }
The official class file \texttt{svjour3.cls} was not found on this machine, so the
document uses a generic article layout. Put the class file into
\texttt{YangTex/templates/classes/} (or install it with \texttt{tlmgr}),
then recompile to obtain the real journal layout.}}
\bigskip
\fi
\begin{abstract}
We establish a Freidlin-Wentzell type large deviation principle for solutions of the two-dimensional stochastic Navier-Stokes equations with small multiplicative noise.
\end{abstract}
\paragraph{Keywords} large deviations; stochastic Navier-Stokes; variational representation; rate function

\section{Introduction}
Consider the stochastic Navier-Stokes system on the two-dimensional torus $\mathbb{T}^2$ driven by a small noise $\sqrt{\eps}\,G(u)\,\dd W$.
For the classical background and related results, see~\cite{Evans,Brezis}.

\begin{theorem}
As $\eps\to0$, the family $\{u^\eps\}$ satisfies a large deviation principle in $C([0,T];H)$ with rate function $I(f)=\frac12\inf\{\|h\|_{L^2}^2:\ f=\mathcal{G}(h)\}$.
\end{theorem}

\begin{proof}
The weak convergence approach reduces the problem to compactness of the controlled skeleton equation and a variational representation.
\end{proof}

\section{Main estimates}
\begin{equation}
\dd u^\eps +\bigl(\nu A u^\eps+B(u^\eps)\bigr)\dd t =\sqrt{\eps}\,G(u^\eps)\,\dd W,
\end{equation}
The rate function is shown to be lower semicontinuous with compact level sets.

\section{Concluding remarks}
The method extends to more general domains and boundary conditions; a detailed
treatment will appear elsewhere.

\begin{thebibliography}{9}
\bibitem{Evans} L.~C. Evans, \emph{Partial Differential Equations}, 2nd ed., AMS, 2010.
\bibitem{Brezis} H. Brezis, \emph{Functional Analysis, Sobolev Spaces and PDEs}, Springer, 2011.
\end{thebibliography}
\end{document}
